% ==========================================================
% Shared preamble: packages, page geometry, custom environments
% ==========================================================
\usepackage[a4paper,margin=1in]{geometry}
\usepackage[utf8]{inputenc}
\usepackage[T1]{fontenc}
\usepackage{lmodern}
\usepackage{amsmath,amssymb,amsthm}
\usepackage{enumitem}
\usepackage{xcolor}
\usepackage{graphicx}
\usepackage{fancyhdr}
\usepackage{titlesec}
\usepackage{tcolorbox}
\tcbuselibrary{skins,breakable}
\usepackage{algorithm}
\usepackage{algpseudocode}
\usepackage{listings}
\usepackage{hyperref}
\usepackage{mdframed}
\usepackage{array}
\usepackage{booktabs}
\usepackage{longtable}
\usepackage{tikz}
\usetikzlibrary{arrows.meta,positioning,decorations.pathreplacing,calc}
\tikzset{every node/.append style={align=center}}

% ---------------- Header / Footer ----------------
\pagestyle{fancy}
\fancyhf{}
\fancyhead[L]{\footnotesize Maruri Saikiran}
\fancyhead[R]{\footnotesize ICPC Training: Greedy, Bitwise, Number Theory}
\fancyfoot[C]{\thepage}
\renewcommand{\headrulewidth}{0.4pt}
\renewcommand{\footrulewidth}{0pt}

% ---------------- Chapter style with epigraph ----------------
\newenvironment{bookepigraph}
  {\list{}{\leftmargin=1.2cm\rightmargin=1.2cm}\item[]\itshape}
  {\endlist}
\newcommand{\chapterepigraph}[1]{%
  \begin{bookepigraph}#1\end{bookepigraph}
}

\titleformat{\chapter}[display]
  {\normalfont\bfseries\Huge}
  {\chaptertitlename\ \thechapter}
  {10pt}{\Huge}
\titlespacing*{\chapter}{0pt}{0pt}{20pt}

% ---------------- Code listing style (C++) ----------------
\definecolor{codebg}{rgb}{0.97,0.97,0.97}
\definecolor{codekw}{rgb}{0.0,0.0,0.6}
\definecolor{codecomment}{rgb}{0.25,0.5,0.25}
\definecolor{codestring}{rgb}{0.6,0.1,0.1}

\lstdefinestyle{cppstyle}{
  language=C++,
  backgroundcolor=\color{codebg},
  basicstyle=\ttfamily\small,
  keywordstyle=\color{codekw}\bfseries,
  commentstyle=\color{codecomment}\itshape,
  stringstyle=\color{codestring},
  numbers=left,
  numberstyle=\tiny\color{gray},
  numbersep=8pt,
  frame=single,
  rulecolor=\color{gray!40},
  breaklines=true,
  showstringspaces=false,
  tabsize=2,
  captionpos=b,
  columns=fullflexible,
  keepspaces=true
}
\lstset{style=cppstyle}

% ---------------- Custom pedagogical boxes ----------------

% Box 1: Pattern Recognition -- keywords that hint at the technique
\newtcolorbox{patternbox}[1][Pattern Recognition]{
  colback=blue!4!white,
  colframe=blue!40!black,
  fonttitle=\bfseries,
  title={#1},
  breakable,
  boxrule=0.6pt,
  arc=2pt,
  left=6pt,right=6pt,top=4pt,bottom=4pt
}

% Box 2: Key Idea / Greedy Choice / Observation
\newtcolorbox{ideabox}[1][Key Idea]{
  colback=yellow!8!white,
  colframe=orange!70!black,
  fonttitle=\bfseries,
  title={#1},
  breakable,
  boxrule=0.6pt,
  arc=2pt,
  left=6pt,right=6pt,top=4pt,bottom=4pt
}

% Box 3: Complexity Analysis
\newtcolorbox{complexitybox}[1][Complexity Analysis]{
  colback=green!5!white,
  colframe=green!35!black,
  fonttitle=\bfseries,
  title={#1},
  breakable,
  boxrule=0.6pt,
  arc=2pt,
  left=6pt,right=6pt,top=4pt,bottom=4pt
}

% Box 4: Key Takeaway / Principle (end of section)
\newtcolorbox{takeawaybox}[1][Key Takeaway]{
  colback=gray!6!white,
  colframe=black!70!white,
  fonttitle=\bfseries,
  title={#1},
  breakable,
  boxrule=0.6pt,
  arc=2pt,
  left=6pt,right=6pt,top=4pt,bottom=4pt
}

% Problem statement leading paragraph
\newcommand{\problemstatement}[1]{\noindent\textbf{Problem Statement.} #1\par}

% Dry run example label
\newcommand{\dryrun}[1]{\noindent\textbf{Dry Run.} #1\par}

\setlist[itemize]{leftmargin=1.4em}
\setlist[enumerate]{leftmargin=1.6em}

% ---------------- Additions for the ICPC training book ----------------

% Box 5: Common pitfalls / edge cases
\newtcolorbox{pitfallbox}[1][Common Pitfalls and Edge Cases]{
  colback=red!4!white,
  colframe=red!55!black,
  fonttitle=\bfseries,
  title={#1},
  breakable,
  boxrule=0.6pt,
  arc=2pt,
  left=6pt,right=6pt,top=4pt,bottom=4pt
}

% Box 6: warning for statements reconstructed from editorials
\newtcolorbox{reconbox}[1][Reconstructed Statement]{
  colback=violet!5!white,
  colframe=violet!60!black,
  fonttitle=\bfseries,
  title={#1},
  breakable,
  boxrule=0.6pt,
  arc=2pt,
  left=6pt,right=6pt,top=4pt,bottom=4pt
}

% Source / difficulty line under a problem title
\newcommand{\source}[2]{\noindent{\small\textbf{Source:} #1 \hfill \textbf{Level:} #2}\par\smallskip}

% Pull a tested solution straight from the repository
\newcommand{\cppfile}[1]{\lstinputlisting{#1}}

\lstset{literate={~}{{\textasciitilde}}1}

% Clickable link to the original problem, shown right under each title
\hypersetup{colorlinks=true,linkcolor=black,urlcolor=blue!70!black}
\newcommand{\problemlink}[2]{\noindent{\small\textbf{Problem link:}
  \href{#1}{#2\ \ensuremath{\nearrow}}}\par}
